% !TEX root = ../../mixed-sklyanin-cubics.tex
\section{孤立 Poisson 零点と被約スキームの不一致}
\label{mix:sec:poisson-limits}

ここでは滑らかな三次超曲面上の\emph{全ての} Poisson 構造を扱う。
入力は \cite[定理\texttt{thm:pg-normal}]{NCCubicA} の全次元正規形であり、
一般三次、一般 Poisson 構造、あるいは各ブロックの楕円性を仮定しない。
以下の結論はこの正規形から直接得られ、比較核の存在や
categorical Torelli を入力としない。

退化軌跡の次元に関する既存結果との違いを確認しておく。
Gualtieri--Pym
\cite[Conjecture~1.1, Theorem~7.13]{GualtieriPymDegeneracy}と、
Druel--Pereira--Pym--Touzet
\cite[Theorem~1.2, Corollary~1.3]{DPPTBondalFive}の結果は、
十分大きな次元の既約成分の存在を保証する。
特に後者は五次元 Fano の Poisson 零集合に正次元成分を与えるが、
孤立成分の併存までは排除しない。
以下の三次超曲面に対する全孤立成分の排除とは、結論と入力が異なる。
この一般次元制約についての新規性の網羅的な確認は未了である。

\begin{theorem}[孤立 Poisson 零点の必要次元]
\label{ppz:thm:isolated-zero-dimension}
$X\subset\mathbb P^{n+1}_{\mathbb C}$ を滑らかな三次超曲面、
$n\ge3$ とし、$\pi\in H^0(X,\bigwedge^2T_X)$ を Poisson 構造とする。
零スキーム $Z_X(\pi)$ の台が孤立点を持つなら、
\[
 n=3k\quad(k\ge1),\qquad \operatorname{rank}_{\mathrm{gen}}\pi=2k.
\]
さらに各孤立点における零スキームは零次元局所完全交叉である。
特に $3\nmid n$ なら、$Z_X(\pi)_{\mathrm{red}}$ は孤立点を持たない。
\end{theorem}

\begin{proof}
$\pi=0$ なら零スキームは $X$ 全体であり、結論の仮定を満たさない。
$\pi\ne0$ とする。引用した正規形により、
$n+2=3k+m$、$k\ge1$、$m\ge0$ となる座標
$x_s=(x_{s1},x_{s2},x_{s3})$、$u=(u_1,\ldots,u_m)$ を取れ、
\[
 F=\sum_{s=1}^k f_s(x_s,u),\qquad
 q=\sum_{s=1}^k q_s,\qquad
 q_s=\iota_{d_{x_s}f_s}
 (\partial_{x_{s1}}\wedge\partial_{x_{s2}}\wedge\partial_{x_{s3}})
\]
となる。$q$ の射影化が $\pi$ で、一般階数は $2k$ である。
$f_s$ は斉次三次式であり、中心変数のみの項はいずれかの $f_s$ に配分する。
$g_s(x_s)=f_s(x_s,0)$ と置く。
この $g_s$ は特異でも零式でもよい。

Euler ベクトル場を
$E=\sum_{s,i}x_{si}\partial_{x_{si}}+\sum_\alpha u_\alpha\partial_{u_\alpha}$
と書く。原点を除く錐では $\mathbb CE$ が接束の直和因子となるので、
零スキームの理想は $E\wedge q$ の係数を $X$ へ制限したものである。
$X$ の滑らかさにより $\bigwedge^2T_X\hookrightarrow
\bigwedge^2T_{\mathbb P^{n+1}}|_X$ も局所的に直和因子であり、
この係数理想は $X$ 上の切断の零理想そのものである。

まず $u_\alpha\ne0$ の chart を考える。
$q$ は中心方向を含まないので、$E\wedge q$ の
$\partial_{u_\alpha}\wedge\partial_{x_{si}}\wedge\partial_{x_{sj}}$
の係数は、符号を除き
$u_\alpha\partial f_s/\partial x_{s\ell}$ である。
従ってこの chart の零理想は、全 $3k$ 個の active 偏微分
\[
 \left(\frac{\partial f_s}{\partial x_{si}}\ \middle|\
       1\le s\le k,\ 1\le i\le3\right)
\]
で生成される。
正則な局所環 $\mathcal O_{X,p}$ は次元 $n$ なので、
この零スキームの各局所成分は
\[
 \dim_p Z_X(\pi)\ge n-3k=m-2
\]
を満たす。孤立点があるなら $m\le2$ である。
一方、この chart があるので $m\ge1$。
$m=1$ の場合、零点では全 active 偏微分が零となり、
$F=0$ と Euler の等式から $u_1F_{u_1}=0$ も得る。
これは $X$ の滑らかさに反する。
従って $u\ne0$ に孤立点があれば $m=2$、すなわち $n=3k$ である。
この場合、その点の零理想は次元 $n$ の正則局所環の
$n=3k$ 個の元で生成される極大イデアル準素な理想である。
従ってこれらの元は正則列となり、零次元局所完全交叉を得る。

残る $u=0$ に孤立点がないことを、$g_s$ の滑らかさを使わずに示す。
点 $p=[x_1:\cdots:x_k:0]$ を取り、$\pi(p)=0$ とする。
もしある $q_s(p)\ne0$ なら、任意の $t\ne s$ について、
$E\wedge q$ のブロック外積次数 $(1\text{ in }t,2\text{ in }s)$
の成分は $x_t\wedge q_s(p)$ である。
その成分に他の項は寄与しないため、$x_t=0$ となる。
従って $p$ は一つの純ブロック平面にあり、
$g_s(x_s)=F(p)=0$ を満たす。
純ブロック平面の閉集合 $V(g_s)\subset\mathbb P^2$
全体が Poisson 零集合に入る。
実際、この平面上では他ブロックの二次偏微分は零であり、
$E\wedge q=3g_s P_s$ となる。
$g_s\ne0$ なら $V(g_s)$ は全ての点で局所次元一以上の平面三次曲線、
$g_s=0$ なら平面全体なので、$p$ は孤立しない。

次に全ての $q_s(p)=0$ の場合を考える。
非零ブロックが一つだけなら前段の純ブロック曲線に属する。
二つ以上あるなら、それぞれの非零 $x_s$ について
$d g_s(x_s)=0$、従って $g_s(x_s)=0$ である。
これらの各ブロックベクトルを独立にスカラー倍して得る射影線型空間は
$X$ と Poisson 零集合に含まれる。
実際、勾配は二次斉次、$g_s$ は三次斉次なので、
その線型空間の全点でも全 active 勾配と $F$ が零となる。
この線型空間は $p$ を通り次元一以上である。
以上により $u=0$ のどの零点も孤立せず、証明が完了する。
\end{proof}

\begin{corollary}[混合五様体による被約対応の反例]
\label{ppz:cor:reduced-point-poisson-counterexample}
例\ref{cmo:ex:sixty-four}の二因子族を取る。
定常ノルム $N(c)$ の六十四個の根のそれぞれにおいて、
$A_c$ は全 coherent 導来圏と全 coherent Kuznetsov 成分が
smooth proper である非可換三次五様体であり、
\[
 \operatorname{Pt}(A_c)=E_1\sqcup E_2\sqcup\{p_c\}
\]
は被約である。
任意の滑らかな可換三次五様体 $Y$ と任意の Poisson 構造 $\pi$ に対して
\[
 \operatorname{Pt}(A_c)_{\mathrm{red}}
   \not\cong Z_Y(\pi)_{\mathrm{red}}.
\]
従って点スキームと Poisson 零スキームの一致という素朴な期待は、
被約スキームの同型へ弱めても一般には成立しない。
\end{corollary}
\begin{proof}
source の smooth proper 性は系\ref{cmo:cor:fixed-zero-norm-smooth}、
点スキームの表示と被約性は例\ref{cmo:ex:sixty-four}による。
左辺には孤立した連結成分 $\{p_c\}$ があるが、
定理\ref{ppz:thm:isolated-zero-dimension}により
$3\nmid5$ の可換三次五様体の Poisson 零集合には孤立点がない。
孤立点の有無は抽象被約スキーム同型で保存されるので結論を得る。
この議論は同じ Ku を持つ標的の存在を仮定せず、
全ての滑らかな可換三次五様体を一括して排除している。
従って主予想の真偽については結論しない。
\end{proof}

\begin{example}[一つの有理係数の反例]
\label{ppz:ex:reduced-counterexample-rational}
$c=-44828$ を取ればよい。
定常 character は $u=1$ において
\[
 (x_1,y_1,z_1)=(-1,-2,-3),\qquad
 (x_2,y_2,z_2)=(-4,-5,-7)
\]
であり、中心値は $28$ と $448$ なので
$28+100\cdot448+c=0$ を満たす。
六十四個のノルム根が相異なることから、この係数で加わる定常点は一つだけである。
$P_{64}$ と $N$ は互いに素なので $P_{64}(-44828)\ne0$ であり、
この例は上の系の smooth proper な source を持つ。
\end{example}

\begin{remark}[既知入力と未解決点]
本節の全 Poisson 構造への適用は、未公表入力\cite{NCCubicA}の
全次元正規形に依存する。正規形の主張だけでなく、
Hessian の基点自由性、\texttt{lem:pg-basepoint-free-lie} の根基排除、
\texttt{lem:pg-quadratic-lie} の三重ベクトル分解を入力箇所として明記する。
その一般的入力と本節の孤立点の排除を区別する。
三様体の非被約局所環を用いる不一致の例は、
可換実現を扱う主稿の系\paperref{main}{ppz:cor:complex-counterexample}で
実際の Ku 同値とともに述べられている。
本節はノルム根での可換 Ku 比較核を構成せず、主予想を否定しない。
この一般次元の制約に関する文献上の新規性は未確認である。
\end{remark}
